% =============================================================================
% AURORA - Informe tecnico: post-procesado QKD en Zynq UltraScale+ (ZCU104)
% Compilar con pdflatex (Overleaf: subir la carpeta docs/report completa).
% Los datos de recursos y potencia salen de datos.tex (tools/report_data.py).
% =============================================================================
\documentclass[a4paper,11pt]{article}
\usepackage[utf8]{inputenc}
\usepackage[T1]{fontenc}
\usepackage{lmodern}
\usepackage[spanish,es-tabla,es-nodecimaldot]{babel}
\usepackage[margin=2.3cm]{geometry}
\usepackage{amsmath}
\usepackage{graphicx}
\usepackage{booktabs}
\usepackage{tabularx}
\usepackage{array}
\usepackage{xcolor}
\usepackage{tikz}
\usetikzlibrary{positioning,fit,arrows.meta,backgrounds}
\usepackage{caption}
\usepackage[hidelinks]{hyperref}

\definecolor{ps}{HTML}{2F6DB5}
\definecolor{pl}{HTML}{C0392B}
\definecolor{ok}{HTML}{1E8449}
\newcommand{\pend}{\textcolor{orange}{\textit{pendiente}}}
\newcommand{\PS}{\textcolor{ps}{\textbf{PS}}}
\newcommand{\PL}{\textcolor{pl}{\textbf{PL}}}
\newcommand{\code}[1]{\texttt{#1}}

% datos.tex esta junto a main.tex. Overleaf compila desde la raiz del proyecto
% si el repositorio se importa entero: se prueban las dos rutas.
\InputIfFileExists{datos}{}{\InputIfFileExists{docs/report/datos}{}{%
  \PackageError{aurora}{No se encuentra datos.tex}{}}}

\title{\textbf{AURORA}\\[2pt]\large Post-procesado QKD (BB84) en un SoC Zynq UltraScale+:\\
codiseño hardware/software, implementación y resultados en placa}
\author{Informe técnico}
\date{Noviembre de 2026}

\begin{document}
\maketitle

\begin{abstract}
\noindent
Se describe la implementación del post-procesado de una distribución cuántica
de claves BB84 en una placa ZCU104 (XCZU7EV). La parte con más volumen de datos
(sincronización, \emph{sifting}, \emph{hash} de Toeplitz) se ejecuta en la
lógica programable (PL) como \emph{kernels} HLS en \emph{streaming}; la
corrección de errores, la autenticación Wegman--Carter del canal clásico y el
control del protocolo se ejecutan en el procesador (PS). Un generador de
números aleatorios verdaderos (TRNG) de osciladores en anillo en la PL
proporciona las semillas públicas, con tests de salud SP~800-90B en hardware.
Sobre 101 capturas reales (QBER 1,72\%), la placa reproduce bit a bit los
resultados del modelo de referencia y entrega \textbf{28\,377 bits de clave
neta} a partir de 40\,755 bits cribados. El TRNG tiene una entropía mínima
medida de 0,785 bits/muestra (NIST \code{ea\_non\_iid}) y supera el test de
reinicios de SP~800-90B. Se incluye, como alternativa a Cascade, una
reconciliación LDPC de un solo sentido.
\end{abstract}

\tableofcontents

% =============================================================================
\section{Plataforma y flujo de diseño}

\begin{itemize}
  \item \textbf{Placa}: Xilinx ZCU104, SoC XCZU7EV-FFVC1156-2-E (4$\times$A53 a 1,2\,GHz + PL UltraScale+).
  \item \textbf{Herramientas}: Vitis HLS, Vivado y Vitis 2022.2; aplicación \emph{bare-metal} (standalone) en el A53\#0.
  \item \textbf{Reloj de la PL}: \code{pl\_clk0} = 100\,MHz (objetivo HLS 10\,ns).
  \item \textbf{Interfaces}: control por \code{M\_AXI\_HPM0\_FPD} (AXI4-Lite). Los \emph{kernels} acceden a la DDR por
        \code{S\_AXI\_HP0\_FPD} a través de una única SmartConnect; el puerto no es coherente, así que el PS
        hace \emph{flush} e \emph{invalidate} de la caché.
  \item \textbf{Flujo reproducible} (\code{Makefile}): \code{make hls-all} $\rightarrow$ \code{make bitstream}
        $\rightarrow$ \code{make baremetal} $\rightarrow$ \code{make board-run}; los mismos fuentes C++ se compilan en el PC
        (tests y modelo de referencia), en HLS (\emph{kernels}) y para el A53.
\end{itemize}

\begin{table}[htbp]
\centering
\caption{Mapa de direcciones (AXI4-Lite).}
\begin{tabular}{lll}
\toprule
Periférico & Base & Implementación \\
\midrule
TRNG & \code{0xA000\_0000} & RTL (Verilog) \\
\code{aurora\_sync} & \code{0xA001\_0000} & HLS \\
\code{aurora\_sift} & \code{0xA002\_0000} & HLS \\
\code{aurora\_toeplitz} & \code{0xA003\_0000} & HLS \\
\bottomrule
\end{tabular}
\end{table}

% =============================================================================
\section{Codiseño hardware/software}

\begin{figure}[h]
\centering
\begin{tikzpicture}[
  blk/.style={draw, rounded corners=2pt, minimum height=8mm, minimum width=27mm, align=center, font=\small},
  psb/.style={blk, fill=ps!12, draw=ps},
  plb/.style={blk, fill=pl!10, draw=pl},
  arr/.style={-{Latex[length=2mm]}, thick},
  node distance=4mm and 6mm]
  % PL
  \node[plb] (sync) {Sync\\\scriptsize histograma de $\Delta t$};
  \node[plb, right=of sync] (sift) {Sift\\\scriptsize coincidencias + bases};
  \node[plb, right=of sift] (toep) {Toeplitz\\\scriptsize PA y extractor};
  \node[plb, right=of toep] (trng) {TRNG\\\scriptsize 4 RO + salud};
  % PS
  \node[psb, below=16mm of sync] (lag) {Lag\\\scriptsize centroide};
  \node[psb, right=of lag] (ec) {Cascade / LDPC\\\scriptsize reconciliación};
  \node[psb, right=of ec] (ver) {Verificación\\\scriptsize hash $2^{61}-1$};
  \node[psb, right=of ver] (auth) {Autenticación\\\scriptsize Wegman--Carter};
  \begin{scope}[on background layer]
    \node[draw=pl, fill=pl!3, rounded corners, inner sep=3mm, fit=(sync)(trng), label={[pl,font=\bfseries\small]above:PL (lógica programable)}] {};
    \node[draw=ps, fill=ps!3, rounded corners, inner sep=3mm, fit=(lag)(auth), label={[ps,font=\bfseries\small]below:PS (A53, bare-metal)}] {};
  \end{scope}
  \draw[arr] (sync) -- (lag);
  \draw[arr] (lag) -- (sift);
  \draw[arr] (sift) -- (ec);
  \draw[arr] (ec) -- (ver);
  \draw[arr] (ver) -- (toep);
  \draw[arr] (trng) -- node[right,font=\scriptsize,align=left]{semillas\\públicas} (auth);
  \node[left=6mm of sync, font=\small, align=center] (in) {Batches\\(API)};
  \draw[arr] (in) -- (sync);
\end{tikzpicture}
\caption{Reparto de la cadena de post-procesado entre la PL y el PS. Las flechas entre bloques indican el orden de la cadena;
todos los datos pasan por la DDR.}
\label{fig:arch}
\end{figure}

\begin{table}[htbp]
\centering
\caption{Reparto de cada etapa y justificación.}
\small
\begin{tabularx}{\textwidth}{l c X}
\toprule
Etapa & Dónde & Motivo \\
\midrule
TRNG (osciladores, tests RCT/APT, FIFO) & \PL & Solo puede ser hardware; tests de salud continuos sin carga para la CPU. \\
Sync: histograma de diferencias de tiempo & \PL & $\sim$52\,000 marcas por batch; ventana deslizante a 1 marca/ciclo. \\
Sync: centroide y elección del lag & \PS & Pocas operaciones sobre 1024 bins. \\
Sifting & \PL & Mezcla de dos flujos ordenados a 1 elemento/ciclo. \\
Cascade (o LDPC) & \PS & Interactivo: dominado por los mensajes con el otro extremo, no por el cálculo. Lógica de control con muchas ramas. \\
Hash de verificación & \PS & Un polinomio sobre $\sim$40\,kbit: despreciable. \\
Amplificación de privacidad (Toeplitz) & \PL & La operación más costosa: $O(n\,m)$ bits. \\
Extractor del TRNG (Toeplitz) & \PL & Mismo \emph{kernel} que la amplificación de privacidad. \\
Autenticación, reserva de clave, protocolo & \PS & Control, red y estado. \\
\bottomrule
\end{tabularx}
\end{table}

\paragraph{¿Cascade en la PL?} No compensa por ahora. Cascade necesita una ida y vuelta por el
canal clásico por cada paso de su búsqueda binaria (cientos por bloque); cuando Alice y Bob estén
en equipos distintos, la latencia de red dominará sobre el cálculo de paridades. Además, el A53 va
sobrado: el sistema genera $\sim$16\,kbit/s de clave cribada y el post-procesado completo en el PS
procesa del orden de $10^5$ bits/s. Si en el futuro la tasa sube a Mbit/s o la latencia del enlace
es alta, la opción natural en la PL es un decodificador LDPC (un solo mensaje por trama, muy
paralelizable); el código LDPC de la sección~\ref{sec:ldpc} está preparado con un decodificador
\emph{min-sum} normalizado, que es el adecuado para hardware.

% =============================================================================
\section{Cadena de post-procesado}

\begin{enumerate}
  \item \textbf{Sincronización} por batch (25\,ms): histograma de $\Delta t = t_A - t_B$ con paso grueso
        (1,024\,ns) y fino (8\,ps); lag por centroide con el fondo restado. Las capturas reales tienen un
        desfase distinto por batch ($\sim$77\,ps rms) y sin deriva dentro del batch: si la sincronización
        falla, el batch se descarta.
  \item \textbf{Sifting}: emparejamiento voraz de detecciones dentro de $\pm400$\,ps y comparación de bases (ventana óptima en el barrido de la sección~\ref{sec:ventana}; los resultados en placa de este informe se midieron con $\pm350$\,ps).
  \item \textbf{Reconciliación}: Cascade (4 pases, $k_1 = 0{,}73/Q$, efecto cascada, límite de trabajo) o
        LDPC (sección~\ref{sec:ldpc}).
  \item \textbf{Verificación}: hash polinómico 2-universal módulo $2^{61}-1$ con semilla nueva del TRNG (64 bits).
  \item \textbf{Amplificación de privacidad}: Toeplitz con semilla completa de $n+m-1$ bits por bloque;
        $m = n\,(1-h(Q)) - \text{fuga}_{EC} - 64 - 2\log_2(1/\varepsilon)$ (fórmula asintótica).
  \item \textbf{Autenticación}: Wegman--Carter, PolyHash sobre $\mathrm{GF}(2^{128})$ cifrado con un OTP de
        128 bits; cada fase se autentica en ambos sentidos (1152 bits de clave por ronda, repuestos con clave QKD).
        La clave se entrega solo si todas las etiquetas verifican.
\end{enumerate}

% =============================================================================
\section{Implementación en la PL}

\subsection{Kernels HLS}
La primera versión de \code{sync} y \code{sift} no cerraba timing (17,2 y 17,6\,ns frente a 10\,ns) porque cada
lectura \code{m\_axi} dependía de la iteración anterior. Se reescribieron como \code{DATAFLOW} en
\emph{streaming}: lectores en ráfaga $\rightarrow$ núcleo con el estado en registros $\rightarrow$ escritor.
En \code{sync}, el núcleo es una ventana deslizante de 32 marcas (máximo observado: 7) con reenvío en el
histograma; si la ventana se desborda, el \emph{kernel} lo indica (\code{complete = 0}) y el PS rechaza el lag.

\begin{table}[htbp]
\centering
\caption{Síntesis HLS (estimación; objetivo 10\,ns; todos los bucles internos con II = 1).
Los recursos tras la implementación son menores (tabla~\ref{tab:modulos}).}
\begin{tabular}{lrrrrrr}
\toprule
Kernel & Estimado (ns) & LUT & FF & BRAM18 & URAM & DSP \\
\midrule
\TablaHLS
\bottomrule
\end{tabular}
\end{table}

\subsection{TRNG}
Cuatro osciladores en anillo de 3, 5, 7 y 9 etapas (LUT con \code{DONT\_TOUCH}) combinados con XOR,
sincronizador de dos biestables y decimación (1 muestra cada 64 ciclos: 1,56\,Mmuestras/s).
Tests de salud SP~800-90B en hardware (RCT, umbral 41; APT, ventana 1024, umbral 793; para
$H_{\min}=0{,}5$ y $\alpha = 2^{-20}$), test de arranque de 1024 muestras, inyección de fallos para el
autodiagnóstico y modo \emph{bypass} solo para evaluación. El extractor (Toeplitz con semilla pública fija,
bloques de 65\,536 bits) se ejecuta en el \emph{kernel} Toeplitz.

\subsection{Implementación (Vivado)}
\begin{itemize}
  \item Timing: \textbf{WNS = +\TimingWNS\,ns}, WHS = +\TimingWHS\,ns, sin \emph{endpoints} que fallen.
  \item DRC: 4 avisos LUTLP-2 (bucles combinacionales permitidos), uno por oscilador en anillo, como se espera.
  \item Potencia estimada: total \PotenciaTotal{} (dinámica \PotenciaDinamica{}, estática \PotenciaEstatica{}), PS incluido.
  \item Kernel LDPC: por iteración, $\sim$($d$+7) ciclos por fila en la comprobación del síndrome y
        $\sim$(2$d$+13) en la actualización ($d$: grado de la fila). Con el código de tasa 0,80
        ($m = 3277$, $d \approx 25$) son $\sim$3,1\,ms por iteración a 100\,MHz.
\end{itemize}

\begin{table}[htbp]
\centering
\small
\begin{minipage}[t]{0.38\textwidth}
\centering
\captionof{table}{Recursos del diseño completo.}
\resizebox{\linewidth}{!}{%
\begin{tabular}{lrrr}
\toprule
Recurso & Usado & Disp. & \% \\
\midrule
\TablaRecursos
\bottomrule
\end{tabular}}
\end{minipage}\hfill
\begin{minipage}[t]{0.58\textwidth}
\centering
\captionof{table}{Recursos por módulo (implementado).}\label{tab:modulos}
\resizebox{\linewidth}{!}{%
\begin{tabular}{lrrrrr}
\toprule
Módulo & LUT & FF & BRAM36/18 & URAM & DSP \\
\midrule
\TablaModulos
\bottomrule
\end{tabular}}
\end{minipage}
\end{table}

% =============================================================================
\section{Resultados en placa}

\subsection{Cadena completa (bare-metal, ZCU104)}
Datos: 101 capturas reales de 25\,ms, 79\,532 coincidencias, 40\,755 bits cribados, QBER 1,718\%.

\begin{table}[htbp]
\centering
\caption{Prueba en placa: resultado \textcolor{ok}{\textbf{TODO OK (0 fallos)}}.}
\small
\begin{tabularx}{\textwidth}{l X}
\toprule
Prueba & Resultado \\
\midrule
TRNG en bruto & 1\,Mbit en 671\,ms (1,56\,Mmuestras/s), sin pérdidas ni alarmas \\
TRNG autodiagnóstico & Los 3 fallos inyectados disparan alarma; arranque correcto \\
TRNG extraído & 100\,000 bits en 615\,ms (162\,kbit/s), equilibrados ($z = 0{,}31$) \\
Sync + sift (3 batches) & Lag, coincidencias y claves \textbf{idénticos bit a bit} al PC; 5,1\,ms + 2,7\,ms por batch de 25\,ms (incluye el trabajo del ARM) \\
Toeplitz & Idéntico a la definición (2000 $\rightarrow$ 1300 bits en 0,34\,ms) \\
Cadena autenticada & 40\,755 cribados $\rightarrow$ 29\,529 tras PA $\rightarrow$ \textbf{28\,377 bits netos}; claves iguales; 709\,ms
  (de ellos, $\sim$430\,ms generando 70\,409 bits de semillas con el TRNG, estimado) \\
\bottomrule
\end{tabularx}
\end{table}

\noindent
Sync + sift ocupan $\sim$7,8\,ms por batch de 25\,ms (31\% del tiempo real), con margen para procesar en
línea. El cuello de botella del post-procesado es la generación de semillas; se puede sacar del camino
crítico generándolas en segundo plano mientras llegan los batches.

\subsection{Ventana de coincidencia}
\label{sec:ventana}
Barrido del semiancho de la ventana sobre los 101 batches (clave neta con Cascade y PA asintótica). Una
ventana mayor recupera más pares, pero también más accidentales (QBER mayor). El óptimo está en
$\pm400$\,ps, que es el valor por defecto desde el 14/11/2026; la ventana de $\pm250$\,ps del algoritmo de
correlación recibido pierde un 13\% de clave.

\begin{table}[htbp]
\centering
\caption{Clave neta frente al semiancho de la ventana de coincidencia.}
\small
\begin{tabular}{lrrrrrrrrr}
\toprule
$\pm w$ (ps) & 150 & 200 & 250 & 300 & 350 & \textbf{400} & 450 & 500 & 600 \\
\midrule
Cribados & 22\,935 & 29\,288 & 34\,485 & 38\,338 & 40\,755 & \textbf{42\,142} & 42\,802 & 43\,116 & 43\,412 \\
QBER (\%) & 1,25 & 1,37 & 1,44 & 1,56 & 1,72 & \textbf{1,87} & 2,04 & 2,24 & 2,56 \\
Neta & 16\,678 & 21\,256 & 24\,945 & 27\,315 & 28\,381 & \textbf{28\,686} & 28\,339 & 27\,562 & 26\,197 \\
\bottomrule
\end{tabular}
\end{table}

\subsection{Entropía del TRNG}



\begin{itemize}
  \item La salida en bruto tiene un sesgo del 2\% ($P(0) = 0{,}521$) y falla los tests SP~800-22, como es de
        esperar; \code{serial\_2}, que cancela el sesgo de primer orden, pasa (no hay correlación apreciable).
  \item La salida \textbf{extraída} (522\,240 bits) pasa los 9 tests SP~800-22 ($p$ entre 0,28 y 0,92).
  \item Se mantiene $H_{\min} = 0{,}5$ en el diseño: margen de 1,57 veces sobre el peor estimador.
\end{itemize}

\begin{table}[htbp]
\centering
\caption{Evaluación SP~800-90B (\code{ea\_non\_iid}, 1\,Mbit en bruto, temperatura ambiente).}
\begin{tabular}{lr}
\toprule
Estimador & $H$ (bit/muestra) \\
\midrule
Most Common Value & 0,937 \\
Collision & 0,910 \\
Markov & 0,940 \\
\textbf{Compression} & \textbf{0,785} \\
t-Tuple / LRS & 0,912 / 0,970 \\
MultiMCW / Lag / MultiMMC / LZ78Y & 0,944 / 0,992 / 0,938 / 0,938 \\
\midrule
$H_{\text{original}}$ & \textbf{0,785} \\
\bottomrule
\end{tabular}
\end{table}

\begin{table}[htbp]
\centering
\caption{Test de reinicios SP~800-90B §3.1.4 (1000 reinicios $\times$ 1000 muestras, \code{ea\_restart}, $H_I = 0{,}785$).
Cordura: frecuencia máxima del valor más común en filas y columnas frente al umbral.}
\small
\begin{tabular}{lccccc}
\toprule
Modo & Descarte & Cordura & $H_r$ & $H_c$ & $\min(H_r, H_c, H_I)$ \\
\midrule
Parada/arranque & 0 & 564 / 648 & 0,779 & 0,803 & \textbf{0,779} \\
Parada/arranque & 1024 & 567 / 648 & 0,781 & 0,778 & \textbf{0,778} \\
Corte de alimentación & 0 & 564 / 648 & 0,810 & 0,795 & \textbf{0,785} \\
\bottomrule
\end{tabular}
\end{table}

\begin{table}[htbp]
\centering
\caption{Entropía mínima ($H_{\text{original}}$, \code{ea\_non\_iid}) frente a la temperatura del chip.}
\begin{tabular}{lccc}
\toprule
Captura & $T$ PL (°C) & Compression & $H_{\text{original}}$ \\
\midrule
Ambiente (sin sensor) & -- & 0,785 & 0,785 \\
Frío (recién encendida) & 28,8 & 0,764 & 0,764 \\
Normal (ventilador, $>$15\,min) & 29,8--30,1 & 0,786 & 0,786 \\
\midrule
\multicolumn{4}{l}{\emph{Bitstream con calefactor (ubicación distinta del TRNG):}} \\
Referencia, calefactor apagado & 24,3 & 0,829 & 0,805 \\
Caliente, recién apagado el calefactor & 60,5 $\rightarrow$ 52,8 & 0,788 & 0,788 \\
Caliente, calefactor encendido & 60,4 & 0,791 & 0,791 \\
\bottomrule
\end{tabular}
\end{table}

\noindent
No hay transitorio de arranque: desde la primera muestra tras habilitar los osciladores, el valor más común
entre reinicios aparece con frecuencia 0,50--0,56, también tras un corte de alimentación real (1000 ciclos
con un relé USB, $\sim$11\,h). La entropía tras reinicio es igual o mayor que la de funcionamiento continuo:
el TRNG queda validado frente a reinicios. Para la prueba en caliente sin acceso físico a la placa se añadió un calefactor en la propia PL
(32 bancos de 512 osciladores en anillo, con \emph{watchdog} en hardware y regulación por software con
límite de temperatura). Llevó la PL de 27,5 a 60\,°C en $\sim$80\,s y la mantuvo estable con 25 de 32
bancos. Se capturó con el calefactor encendido y justo después de apagarlo (el chip se enfría
rápidamente: de 60,5 a 52,8\,°C durante los 0,7\,s de la captura). Entre 24 y 60\,°C la entropía se
mantiene en 0,76--0,81 sin tendencia con la temperatura, y el ruido de alimentación del calefactor no la
afecta. El margen respecto a $H_{\min}=0{,}5$ de diseño es de 1,5 veces o más en todos los casos.

% =============================================================================
\section{Reconciliación: Cascade frente a LDPC}
\label{sec:ldpc}

La reconciliación LDPC es de un solo sentido por trama: Alice envía el síndrome y Bob responde con éxito y
número de errores. Se usa una familia de códigos IRA ($n = 16\,384$, tasas 0,50--0,95) construidos de forma
determinista (sin ciclos de longitud 4 hasta $R = 0{,}90$), adaptación de tasa con perforado
\emph{untainted} y acortado, QBER adaptativo trama a trama, y decodificación BP por capas
(\emph{sum-product} o \emph{min-sum} normalizado). Las tramas que fallan se descartan.

\begin{table}[htbp]
\centering
\caption{Tasa de error de trama (FER) del LDPC con $Q = 1{,}72\%$ (simulación, sum-product, 100 iteraciones).}
\begin{tabular}{lccccccc}
\toprule
$f$ objetivo & 1,20 & 1,25 & 1,30 & 1,35 & 1,40 & 1,45 & 1,50 \\
\midrule
FER & 13\% & 8\% & 8\% & 2,5\% & 0/124 & 0/120 & 0/120 \\
\bottomrule
\end{tabular}
\end{table}

\begin{table}[htbp]
\centering
\caption{Cadena autenticada sobre los 40\,755 bits cribados reales (PC).}
\resizebox{\textwidth}{!}{%
\begin{tabular}{lcccccc}
\toprule
Método & QBER estimado & $f$ & Fuga & Clave neta & Mensajes & Tiempo (PC) \\
\midrule
Cascade & 2,00\% & 1,172 & 5\,985 & 28\,381 & $\sim$6\,000 consultas & 0,05\,s \\
Cascade & 1,72\% & 1,154 & 5\,891 & 28\,475 & $\sim$6\,000 consultas & 0,05\,s \\
LDPC    & 2,00\% & 1,480 & 7\,556 & 26\,810 & 3 + 3 & 0,96\,s \\
LDPC    & 1,72\% & 1,396 & 7\,125 & 27\,241 & 3 + 3 & 0,68\,s \\
\bottomrule
\end{tabular}}
\end{table}

\subsection{Decodificador LDPC en la PL}
El \emph{kernel} \code{aurora\_ldpc\_dec} implementa la decodificación de Bob con \emph{min-sum} normalizado
por capas en coma fija: LLR de entrada y mensajes de 8 bits, posteriores de 16 bits y normalización
$\alpha = 0{,}6875 = \tfrac12+\tfrac14-\tfrac1{16}$ (solo sumas y desplazamientos). El código (matriz $H$
en CSR) se guarda en memoria interna (URAM/BRAM) y solo se recarga al cambiar de código. El mismo fuente
C++ es el modelo de software, de modo que la comparación PL/PS es exacta bit a bit. Las etapas de
planificación, perforado, síndrome de Alice y mensajes siguen en el PS.

\begin{table}[htbp]
\centering
\caption{FER del decodificador en coma fija ($Q = 1{,}72\%$, 300 tramas por punto) frente a sum-product en flotante.}
\begin{tabular}{lcccc}
\toprule
& \multicolumn{3}{c}{Min-sum coma fija (PL)} & Sum-product (PS) \\
$f$ objetivo & 1,50 & 1,55 & 1,60 & 1,40 \\
\midrule
FER & 1,7\% & 0,7\% & 1,0\% & 0/151 \\
\bottomrule
\end{tabular}
\end{table}

\noindent
El \emph{min-sum} en coma fija necesita más margen ($f = 1{,}50$) y tiene un suelo de error de $\sim$1\%:
con los datos reales entrega $\sim$25\,800 bits netos frente a 27\,241 con sum-product. A cambio, es la
variante que se puede llevar a hardware.

\begin{table}[htbp]
\centering
\caption{Decodificación LDPC de los 40\,755 bits reales en la placa: PL frente a PS (\code{make board-run}, sección 6).}
\begin{tabular}{lccc}
\toprule
Decodificador & Tiempo & Resultado & Clave neta \\
\midrule
PL, min-sum coma fija & \textbf{384\,ms} & 4 tramas, 1 fallida & 25\,845 \\
PS (A53), min-sum coma fija & 502\,ms & 4 tramas, 1 fallida & 25\,841 \\
PS (A53), sum-product flotante & 617\,ms & 3 tramas, 0 fallidas & -- \\
\bottomrule
\end{tabular}
\end{table}

\noindent
La PL y el A53 obtienen \textbf{resultados idénticos bit a bit} (mismas claves, fuga, tramas fallidas e
iteraciones), y la PL es \textbf{1,3 veces más rápida} con el mismo algoritmo, en línea con la estimación por
síntesis ($\sim$3,1\,ms por iteración frente a $\sim$3,9\,ms en el A53). La clave neta difiere en unos pocos
bits entre ejecuciones porque las semillas de la amplificación de privacidad salen del TRNG.

\noindent
Medido en la placa: el A53 reconcilia los 40\,755 bits en 0,5--0,6\,s, mientras que se tardan $\sim$2,5\,s
en generarlos, así que el PS ya sigue el ritmo con un margen de $\sim$4 veces (la estimación previa de
2--3\,s era pesimista). La construcción de los códigos tarda 3,2\,s, una sola vez al arrancar. En la
ejecución en coma fija, la trama que falla agota las 100 iteraciones y se lleva buena parte del tiempo.
El \emph{kernel} procesa una arista por ciclo (100\,M aristas/s): con $\sim$83\,000 aristas, tres pasadas
por iteración y $\sim$10 iteraciones, unos 30\,ms por trama, del mismo orden que el A53. Con las tasas
actuales la PL no es necesaria para el LDPC; lo sería si la tasa de clave creciera. Para una ganancia grande
habría que procesar varias filas en paralelo, lo que exige códigos cuasi-cíclicos (QC-LDPC): es el
siguiente paso si la medida en placa lo justifica.

\noindent
Con este tamaño de trama, el LDPC entrega un 4\% menos de clave que Cascade, pero con 6 mensajes en lugar
de miles de consultas: es la opción a elegir cuando la latencia de red domine o cuando se lleve la
decodificación a la PL. La eficiencia mejoraría con tramas mayores ($n \geq 10^5$) y distribuciones de grados
optimizadas por evolución de densidad.

% =============================================================================
\section{Verificación}

\begin{itemize}
  \item \textbf{Tests en PC} (\code{make test-all}): 73 tests C++ (\emph{kernels}, autenticación, LDPC, TRNG, driver),
        34 tests Python (incluida la validación sobre los 101 batches reales contra un modelo independiente),
        tests que documentan los 9 fallos de la implementación anterior, y 6 simulaciones RTL del TRNG (Icarus).
  \item \textbf{C-simulation HLS} de los tres \emph{kernels} con el mismo \emph{testbench}.
  \item \textbf{Plataforma simulada}: la aplicación \emph{bare-metal} completa se ejecuta en el PC con los periféricos
        emulados a nivel de registro.
  \item \textbf{Placa}: resultados idénticos bit a bit a los del PC; evaluación de entropía con las herramientas de NIST.
  \item \textbf{Seguridad del canal}: se comprueba que cualquier alteración de paridades, síndromes, mensajes o
        etiquetas, o la repetición de una etiqueta, aborta la ronda sin entregar clave.
\end{itemize}

% =============================================================================
\section{Limitaciones y trabajo futuro}

\begin{itemize}
  \item \textbf{Seguridad}: la longitud de la amplificación de privacidad usa la fórmula asintótica de BB84, sin
        correcciones de tamaño finito; debe validarse con el equipo de UPM.
  \item \textbf{Detectores desajustados (hallazgo en los datos reales)}: los retardos de los detectores difieren
        hasta $\sim$120\,ps y sus eficiencias hasta un 40\%. En la base H/V, la diferencia de tiempos de cada par
        da 0,16 bits de información sobre el bit (0,03 compensando los retardos de cada canal), y la clave cribada
        está sesgada ($P(0) = 0{,}38$--$0{,}40$). Si las marcas de tiempo se intercambian por un canal público,
        hay que compensar los retardos (idealmente en el etiquetador de tiempos) e incluir la fuga, el sesgo y el
        desajuste de eficiencia en el análisis de seguridad.
  \item \textbf{Integración con la API del receptor} y paso a Linux (PetaLinux) con \emph{buffers} contiguos.
  \item \textbf{Alice y Bob en procesos separados}, con los anuncios del \emph{sifting} dentro de la transcripción autenticada.
  \item \textbf{Rendimiento en el PS}: medir el reparto de tiempos por etapa; generar las semillas en segundo plano;
        representar las claves empaquetadas (64 bits por palabra) y usar NEON (\code{PMULL} para $\mathrm{GF}(2^{128})$).
  \item \textbf{QRNG del receptor}: si la API lo expone, combinarlo con el TRNG (XOR).
  \item \textbf{LDPC}: tramas mayores, distribuciones optimizadas; en la PL, códigos QC-LDPC para decodificar varias filas en paralelo.
  \item \textbf{TRNG}: evaluación en caliente (tapando el ventilador o en cámara climática) cuando haya acceso al laboratorio.
\end{itemize}

\end{document}
